\begin{hint}[Exercise VI.17.1]
Use $X=D(1)$ and $A_1=A$.
\end{hint}

\begin{hint}[Exercise VI.17.2]
$D(x^3)=D(x)$ and localizing at $x^3$ already makes $x$ invertible.
\end{hint}

\begin{hint}[Exercise VI.17.3]
Send $a/x^m$ to the same fraction and then allow powers of $y$ in denominators.
\end{hint}

\begin{hint}[Exercise VI.17.4]
Every denominator outside $\mfp$ becomes invertible in the colimit, and every element of $A_{\mfp}$ has such a denominator.
\end{hint}

\begin{hint}[Exercise VI.17.5]
A fraction $a/s$ is a unit exactly when $a\notin\mfp$.
\end{hint}

\begin{hint}[Exercise VI.17.6]
First quotient by $\mfp$, then observe that all nonzero classes are inverted.
\end{hint}

\begin{hint}[Exercise VI.17.7]
At $(0)$ every nonzero element is inverted.
\end{hint}

\begin{hint}[Exercise VI.17.8]
Localize at the multiplicative set of polynomials not divisible by $x-a$.
\end{hint}

\begin{hint}[Exercise VI.17.9]
The denominator must not be divisible by $p$; then quotient the local ring by $p$.
\end{hint}

\begin{hint}[Exercise VI.17.10]
Invert $30$, equivalently invert $2,3,5$.
\end{hint}

\begin{hint}[Exercise VI.17.11]
An invertible nilpotent forces $1=0$.
\end{hint}

\begin{hint}[Exercise VI.17.12]
Use the criterion for a germ $a/1$ to be a unit in $A_{\mfp}$.
\end{hint}

\begin{hint}[Exercise VI.17.13]
Recall the zero criterion in a localization.
\end{hint}

\begin{hint}[Exercise VI.17.14]
The spaces are the same one point, but compare global section rings.
\end{hint}

\begin{hint}[Exercise VI.17.15]
Use the two primitive idempotents and the two clopen singleton opens.
\end{hint}

\begin{hint}[Exercise VI.17.16]
This is exactly the two-open sheaf equalizer condition.
\end{hint}

\begin{hint}[Exercise VI.17.17]
Here $D(xy)=\varnothing$, so there is no overlap compatibility condition.
\end{hint}

\begin{hint}[Exercise VI.17.18]
Write $\alpha=a/s$ with $s\notin\mfp$ and use $D(s)$.
\end{hint}

\begin{hint}[Exercise VI.17.19]
Equality in $A_{\mfp}$ means some element outside $\mfp$ annihilates the cross-multiplied difference.
\end{hint}

\begin{hint}[Exercise VI.17.20]
Use $\kappa(0)=\mathbb Q$ and $\kappa(p)=\mathbb F_p$.
\end{hint}

\begin{hint}[Exercise VI.17.21]
At $(x)$ the element $y$ becomes invertible and forces $x=0$; similarly for $(y)$.
\end{hint}

\begin{hint}[Exercise VI.17.22]
Localization commutes with quotient.  In particular,\[\bigl(A/\Nil A\bigr)_{\mfp/\Nil A}\cong A_{\mfp}/(\Nil A)_{\mfp}.\]
\end{hint}

\begin{hint}[Exercise VI.17.23]
Apply the sheaf axiom to the cover $X=D(f)\cup D(g)$ and then identify $\OO_X(X)$ with $A$.
\end{hint}

\begin{hint}[Exercise VI.17.24]
Local fraction representatives remain valid after shrinking.
\end{hint}

\begin{hint}[Exercise VI.17.25]
Equality of all germs gives local equality near each point; then apply locality of the sheaf.
\end{hint}

\begin{hint}[Exercise VI.17.26]
Besides the homeomorphism $\Spec A_f\cong D(f)$, identify the sheaves on a basis and check local rings on stalks.
\end{hint}

